\documentclass[11pt]{article}
\usepackage[utf8]{inputenc}
\usepackage[T1]{fontenc}
\usepackage[brazil]{babel}
\usepackage{amsmath,amssymb,amsthm}
\usepackage{tikz-cd}
\usetikzlibrary{babel}

\newtheorem{teorema}{Teorema}
\newtheorem{definicao}{Definição}
\newcommand{\im}{\operatorname{im}}
\newcommand{\Z}{\mathbb{Z}}

\begin{document}

\section{Diagramas comutativos}

\begin{definicao}
Um quadrado de homomorfismos
\[
\begin{tikzcd}
A \arrow[r, "f"] \arrow[d, "\alpha"'] & B \arrow[d, "\beta"] \\
C \arrow[r, "g"'] & D
\end{tikzcd}
\]
\emph{comuta} quando $\beta\circ f = g\circ\alpha$.
\end{definicao}

\begin{teorema}[Primeiro teorema do isomorfismo]
Seja $\varphi\colon G\to H$ um homomorfismo de grupos e $\pi\colon G\to G/\ker\varphi$ a projeção. Existe um único isomorfismo $\bar\varphi$ tal que o diagrama
\[
\begin{tikzcd}
G \arrow[r, "\varphi"] \arrow[d, "\pi"'] & \im\varphi \\
G/\ker\varphi \arrow[ur, "\bar\varphi"', "\cong"]
\end{tikzcd}
\]
comuta.
\end{teorema}

\begin{proof}
Defina $\bar\varphi(g\ker\varphi)=\varphi(g)$. Está bem definido, pois $g\ker\varphi=h\ker\varphi$ implica $h^{-1}g\in\ker\varphi$, logo $\varphi(g)=\varphi(h)$. É sobrejetor sobre $\im\varphi$ por construção e injetor porque $\varphi(g)=e$ implica $g\in\ker\varphi$.
\end{proof}

\section{Sequências exatas}

Um exemplo clássico com duas linhas:
\[
\begin{tikzcd}
0 \arrow[r] & \Z \arrow[r, "\cdot 2"] \arrow[d, "\cdot 3"'] & \Z \arrow[r] \arrow[d, "\cdot 3"] & \Z/2\Z \arrow[r] \arrow[d, "\cdot 3"] & 0 \\
0 \arrow[r] & \Z \arrow[r, "\cdot 2"'] & \Z \arrow[r] & \Z/2\Z \arrow[r] & 0
\end{tikzcd}
\]
Cada quadrado comuta, porque multiplicações por inteiros comutam.

\end{document}
